\section*{अध्याय \ref{ch:g11:redox} --- ऑक्सीकरण और अपचयन}

\begin{solution}{exo:g11:redox:1}
\omterm{def:g11:redox:oxidant}{ऑक्सीकारक}: वह स्पीशीज़ जो \omterm{def:g9:inside-the-atom:nucleus}{इलेक्ट्रॉन} ग्रहण कर सकती है। \omterm{def:g11:redox:oxidant}{अपचायक}: वह स्पीशीज़ जो \omterm{def:g9:inside-the-atom:nucleus}{इलेक्ट्रॉन}
दे सकती है। \omterm{def:g11:redox:oxidation}{ऑक्सीकरण}: \omterm{def:g9:inside-the-atom:nucleus}{इलेक्ट्रॉनों} का खोना। \omterm{def:g11:redox:oxidation}{अपचयन}: \omterm{def:g9:inside-the-atom:nucleus}{इलेक्ट्रॉनों}
का ग्रहण।
\end{solution}

\begin{solution}{exo:g11:redox:2}
\ce{Fe^{2+}}/\ce{Fe}; \ce{Ag+}/\ce{Ag}; \ce{I2}/\ce{I-};
\ce{Fe^{3+}}/\ce{Fe^{2+}}।
\end{solution}

\begin{solution}{exo:g11:redox:3}
\ce{Ag+ + e- <=> Ag}; \ce{Al^{3+} + 3e- <=> Al};
\ce{Fe^{3+} + e- <=> Fe^{2+}}; \ce{Cl2 + 2e- <=> 2Cl-}।
\end{solution}

\begin{solution}{exo:g11:redox:4}
\omterm{def:g11:redox:oxidation}{ऑक्सीकरण} (\omterm{def:g9:inside-the-atom:nucleus}{इलेक्ट्रॉन} खोए जाते हैं); लोहा \omterm{def:g11:redox:oxidant}{अपचायक} की तरह काम करता है।
\end{solution}

\begin{solution}{exo:g11:redox:5}
जस्ता ऑक्सीकृत होता है और \omterm{def:g11:redox:oxidant}{अपचायक} है; \ce{Cu^{2+}} अपचयित होता है और
\omterm{def:g11:redox:oxidant}{ऑक्सीकारक} है। प्रति जस्ता \omterm{def:g7:atoms-and-molecules:atom}{परमाणु} दो \omterm{def:g9:inside-the-atom:nucleus}{इलेक्ट्रॉन}।
\end{solution}

\begin{solution}{exo:g11:redox:6}
\ce{H2O2 -> H2O} से शुरू कीजिए: % ce-unbalanced-ok (step 1 of the method)
ऑक्सीजन: बाईं ओर दो O, इसलिए दाईं ओर \ce{2H2O}; हाइड्रोजन: बाईं ओर 2,
दाईं ओर 4, इसलिए बाईं ओर \ce{2H+}; आवेश: बाईं ओर $+2$,
दाईं ओर 0, इसलिए बाईं ओर दो \omterm{def:g9:inside-the-atom:nucleus}{इलेक्ट्रॉन}:
\ce{H2O2 + 2H+ + 2e- <=> 2H2O}।
\end{solution}

\begin{solution}{exo:g11:redox:7}
\ce{I2 + 2e- -> 2I-} और \ce{2S2O3^{2-} -> S4O6^{2-} + 2e-}; दोनों में दो-दो
\omterm{def:g9:inside-the-atom:nucleus}{इलेक्ट्रॉन}, इसलिए \ce{I2 + 2S2O3^{2-} -> 2I- + S4O6^{2-}}।
\end{solution}

\begin{solution}{exo:g11:redox:8}
\ce{Cu -> Cu^{2+} + 2e-} और $2\times$ (\ce{Ag+ + e- -> Ag}):
\ce{Cu + 2Ag+ -> Cu^{2+} + 2Ag}। बनने वाले \ce{Cu^{2+}} \omterm{def:g9:ions:ion}{आयन} नीले होते हैं।
\end{solution}

\begin{solution}{exo:g11:redox:9}
\ce{C6H8O6 -> C6H6O6 + 2H+ + 2e-} और \ce{I2 + 2e- -> 2I-}:
\ce{C6H8O6 + I2 -> C6H6O6 + 2I- + 2H+}। विटामिन C \omterm{def:g11:redox:oxidant}{अपचायक} है।
\end{solution}

\begin{solution}{exo:g11:redox:10}
\ce{Al -> Al^{3+} + 3e-} 3 \omterm{def:g9:inside-the-atom:nucleus}{इलेक्ट्रॉन} देता है, \ce{Cu^{2+} + 2e- -> Cu}
2 लेता है; सबसे छोटी उभयनिष्ठ संख्या 6 है, इसलिए पहले को
2 से और दूसरे को 3 से गुणा किया जाता है: \ce{2Al + 3Cu^{2+} -> 2Al^{3+} + 3Cu}।
\end{solution}

\begin{solution}{exo:g11:redox:11}
\ce{H2O2 + 2H+ + 2e- -> 2H2O} और $2\times$ (\ce{Fe^{2+} -> Fe^{3+} + e-}):
\ce{H2O2 + 2H+ + 2Fe^{2+} -> 2H2O + 2Fe^{3+}}।
\end{solution}

\begin{solution}{exo:g11:redox:12}
$n(\ce{Zn}) = 1.30/65.4 = \qty{1.99e-2}{mol}$। ऑक्सीकृत होने वाले हर जस्ता \omterm{def:g7:atoms-and-molecules:atom}{परमाणु}
पर ताँबे का एक \omterm{def:g7:atoms-and-molecules:atom}{परमाणु} जमता है, इसलिए $n(\ce{Cu}) = \qty{1.99e-2}{mol}$
और $m(\ce{Cu}) = 1.99\times 10^{-2} \times 63.5 = \qty{1.26}{g}$।
\end{solution}

\begin{solution}{exo:g11:redox:13}
\ce{Cr2O7^{2-} + 14H+ + 6e- -> 2Cr^{3+} + 7H2O} और
$6\times$ (\ce{Fe^{2+} -> Fe^{3+} + e-}):
\ce{Cr2O7^{2-} + 14H+ + 6Fe^{2+} -> 2Cr^{3+} + 6Fe^{3+} + 7H2O}।
प्रति डाइक्रोमेट \omterm{def:g9:ions:ion}{आयन} छह \ce{Fe^{2+}}: $6 \times 0.010 = \qty{0.060}{mol}$।
\end{solution}

\begin{solution}{exo:g11:redox:14}
\ce{2ClO- + 4H+ + 2e- -> Cl2 + 2H2O} और \ce{2Cl- -> Cl2 + 2e-}; जोड़कर
और 2 से भाग देकर: \ce{ClO- + Cl- + 2H+ -> Cl2 + H2O}। अम्ल
\ce{H+} \omterm{def:g9:ions:ion}{आयन} देता है, और अभिक्रिया डाइक्लोरीन, एक विषैली गैस, छोड़ती है:
इसीलिए दोनों उत्पादों को कभी न मिलाने की चेतावनी।
\end{solution}

\begin{solution}{exo:g11:redox:15}
\omterm{def:g9:inside-the-atom:nucleus}{इलेक्ट्रॉन} विलयन से होकर मुक्त रूप से नहीं जा सकते
(\cref{prop:g11:redox:no-free-electrons}): जस्ता \omterm{def:g7:atoms-and-molecules:atom}{परमाणुओं} द्वारा खोए गए \omterm{def:g9:inside-the-atom:nucleus}{इलेक्ट्रॉन}
केवल \omterm{def:g9:periodic-table-first-look:metal}{धातु} को छूने वाला \ce{Cu^{2+}} \omterm{def:g9:ions:ion}{आयन} ही ले सकता है। इसलिए ताँबे का
\omterm{def:g7:atoms-and-molecules:atom}{परमाणु} पट्टी की सतह पर बनता है, और वहीं रह जाता है।
\end{solution}

\begin{solution}{pb:g11:redox:1}
\textbf{1.} \ce{Cr2O7^{2-}}/\ce{Cr^{3+}} और \ce{CH3COOH}/\ce{CH3CH2OH}।

\textbf{2.} एथेनॉल ऑक्सीकृत होता है: वह \omterm{def:g11:redox:oxidant}{अपचायक} है।

\textbf{3.} \ce{Cr2O7^{2-} + 14H+ + 6e- <=> 2Cr^{3+} + 7H2O}।

\textbf{4.} कार्बन संतुलित है (हर ओर दो)। ऑक्सीजन: बाईं ओर दो,
दाईं ओर एक, इसलिए दाईं ओर \ce{H2O}। हाइड्रोजन: बाईं ओर चार,
दाईं ओर $6 + 2 = 8$, इसलिए बाईं ओर \ce{4H+}। आवेश: बाईं ओर $+4$,
दाईं ओर 0, इसलिए बाईं ओर चार \omterm{def:g9:inside-the-atom:nucleus}{इलेक्ट्रॉन}:
\ce{CH3COOH + 4H+ + 4e- <=> CH3CH2OH + H2O}।

\textbf{5.} बारह \omterm{def:g9:inside-the-atom:nucleus}{इलेक्ट्रॉन}: पहले का $2\times$, उलटे किए गए दूसरे का
$3\times$:
\ce{2Cr2O7^{2-} + 3CH3CH2OH + 16H+ -> 4Cr^{3+} + 3CH3COOH + 11H2O}।

\textbf{6.} जहाँ भी एथेनॉल पहुँचता है, नारंगी \ce{Cr2O7^{2-}} \omterm{def:g9:ions:ion}{आयन} हरे
\ce{Cr^{3+}} \omterm{def:g9:ions:ion}{आयनों} में अपचयित हो जाते हैं।

\textbf{7.} $M = 2\times 39.1 + 2\times 52.0 + 7\times 16.0 =
\qty{294.2}{g/mol}$।

\textbf{8.} $n = 2.0\times 10^{-3}/294.2 = \qty{6.8e-6}{mol}$।

\textbf{9.} प्रति दो डाइक्रोमेट \omterm{def:g9:ions:ion}{आयन} एथेनॉल के तीन \omterm{def:g7:atoms-and-molecules:molecule}{अणु}:
$n = \tfrac{3}{2}\times 6.8\times 10^{-6} = \qty{1.0e-5}{mol}$।

\textbf{10.} $M(\ce{C2H6O}) = \qty{46.0}{g/mol}$, इसलिए
$m = 1.02\times 10^{-5}\times 46.0 = \qty{4.7e-4}{g}$, लगभग
\qty{0.47}{mg}।

\textbf{11.} $0.25/0.47 \approx 0.53$: डाइक्रोमेट का लगभग
\qty{53}{\percent}।

\textbf{12.} $0.53 \times 10 \approx \qty{5.3}{cm}$।

\textbf{13.} $\qty{0.25}{mg} \times 2100 = \qty{525}{mg}$ प्रति लीटर
रक्त, लगभग \qty{0.53}{g/L}।

\textbf{14.} साँस सबसे पहले प्रवेश-द्वार पर दानों से मिलती है; वहाँ
डाइक्रोमेट अधिकता में है और सारा एथेनॉल हटा देता है, इसलिए
नली के पहले हिस्से में अभिक्रिया पूरी हो जाती है, जो पूरी तरह
हरा हो जाता है, इससे पहले कि कोई एथेनॉल आगे पहुँचे।

\textbf{15.} अभिक्रिया पूर्ण है: हरे \ce{Cr^{3+}} \omterm{def:g9:ions:ion}{आयन} वापस
डाइक्रोमेट नहीं बनते, इसलिए इस्तेमाल की हुई नली दोबारा माप नहीं सकती।

% ledger: ghs:K2Cr2O7 (H350 among its hazard statements)
\textbf{16.} यह तीव्र विषैला है (GHS06) और स्वास्थ्य के लिए गंभीर ख़तरा
(GHS08: इसके ख़तरा-कथनों में ``कैंसर का कारण बन सकता है'' शामिल है), साथ ही
संक्षारक और जलीय जीवन के लिए विषैला; आम लोगों द्वारा इस्तेमाल होने वाले उपकरण में
यह नहीं होना चाहिए।

\textbf{17.} \ce{O2 + 4H+ + 4e- -> 2H2O} और
\ce{CH3CH2OH + H2O -> CH3COOH + 4H+ + 4e-}:
\ce{CH3CH2OH + O2 -> CH3COOH + H2O}।

\textbf{18.} चार \omterm{def:g9:inside-the-atom:nucleus}{इलेक्ट्रॉन}।

\textbf{19.} $n = 0.25\times 10^{-3}/46.0 = \qty{5.4e-6}{mol}$
एथेनॉल; \omterm{def:g9:inside-the-atom:nucleus}{इलेक्ट्रॉन} $4n = \qty{2.2e-5}{mol}$, अर्थात
$2.2\times 10^{-5}\times 6.02\times 10^{23} = \num{1.3e19}$ \omterm{def:g9:inside-the-atom:nucleus}{इलेक्ट्रॉन}।

\textbf{20.} $Q = 1.3\times 10^{19}\times 1.60\times 10^{-19} \approx
\qty{2.1}{C}$। एथेनॉल का हर \omterm{def:g7:atoms-and-molecules:molecule}{अणु} तार से होकर ठीक चार \omterm{def:g9:inside-the-atom:nucleus}{इलेक्ट्रॉन}
भेजता है, इसलिए आवेश एथेनॉल की मात्रा के समानुपाती है:
उसे मापना एथेनॉल को मापना है।
\end{solution}
